%%%%%%%%%%%%%%%%%%%%%%%%%%%%%%%%%%%%%%%%%%%%%%%%%%%%%%%%%%%%%%%%%%%%%%
% Lớp 11 | Đạo hàm | Đạo hàm
% Dạng 1D7?1-5: Toán thực tế khác áp dụng định nghĩa đạo hàm
% Số câu: 11  (sinh tự động bởi tach_cau.py; từ nay thầy sửa trực tiếp file này)
%%%%%%%%%%%%%%%%%%%%%%%%%%%%%%%%%%%%%%%%%%%%%%%%%%%%%%%%%%%%%%%%%%%%%%

%%%%%%%%%%%%%%%% TRẢ LỜI NGẮN %%%%%%%%%%%%%%%%

% nguồn: Ok_TraLoiNgan_Lop11_Lan1
\begin{ex}%[1D7H1-5]%[ma:0d1ab7d0]
	Nhiệt độ cơ thể của một người trong thời gian bị bệnh được cho bởi công thức \break$T(t)=-0{,}1t^2+1{,}2t+98{,}6$, trong đó $T$ là nhiệt độ (tính theo đơn vị đo Fahrenheit) tại thời điểm $t$ (tính  theo ngày). Tìm tốc độ thay đổi nhiệt độ ở thời điểm $t=2$.
	\shortans{0,8} 
	 \loigiai{
		Đạo hàm của hàm số $T$ biểu thị  tốc độ thay đổi của nhiệt độ.\\
		 Ta có $T'(t)=-0{,}2t+1{,}2$.\\
		Vậy tốc độ thay đổi nhiệt độ tại thời  điểm $t=2$ là $T'(2)=-0{,}2\cdot 2+1{,}2=0{,}8^\circ F/$ngày.
	} 
\end{ex}

% nguồn: Ok_TraLoiNgan_Lop11_Lan1
\begin{ex}%[1D7H1-5]%[ma:56d0523c]
	Một người gửi tiết kiệm khoản tiền $ 100 $ triệu đồng vào một ngân hàng với lãi suất $7 \% $/năm. Tổng số tiền vốn và lãi (làm tròn đến hàng đơn vị) mà người đó nhận được sau 1  năm là bao nhiêu triệu đồng, nếu tiền lãi được tính theo hình thức lãi kép với kì hạn 6 tháng.
	\shortans{107} 
	 \loigiai{
		Số tiền vốn và lãi người đó nhận được sau một năm theo hình thức lãi kép với kì hạn 6 tháng  là
		$$ 
		T=A  \cdot\left(1+\dfrac{r}{2}\right)^2
		=100(1+3,5  \%)^2=107.
		 $$
	 }
\end{ex}

% nguồn: Ok_TraLoiNgan_Lop11_Lan1
\begin{ex}%[1D7V1-5]%[ma:b617c3cf]
	Một vật chuyển động theo quy luật $s=\dfrac{1}{3}t^3-t^2+9 t$, với $t$ (giây) là khoảng thời gian tính từ lúc vật bắt  đầu chuyển động và $s$ (mét) là quãng đường vật đi được trong thời gian đó. Hỏi trong khoảng thời gian $10$ giây, kể từ lúc bắt đầu chuyển động, vận tốc nhỏ nhất của vật đạt được bằng bao nhiêu?\\
	 \shortans{8}
	 \loigiai{
		 Vận tốc $v$ của vật được tính theo công thức $v(t)=s^{\prime}(t)=t^2-2 t+9$.\\
		Ta có $t^2-2  t+9=(t-1)^2+8\geq 8\Rightarrow v\geq 8$.\\
	Vậy vận tốc nhỏ nhất của vật là $8$ m/s đạt được tại thời điểm  $t=1$ giây.}
\end{ex}

% nguồn: Ok_TraLoiNgan_Lop11_Lan1
\begin{ex}%[1D7V1-5]%[ma:8d72a454]
	Một chất điểm chuyển động theo quy luật $s(t)=-t^3+6t^2$ với $t$ là thời gian  tính từ lúc bắt đầu chuyển động, $s(t)$ là quãng đường đi được trong khoảng thời gian $t$. Tính thời điểm $t$ tại đó vận tốc đạt giá trị lớn nhất.
	 \shortans{2}
	\loigiai{Điều kiện $t>0$.\\ 
	Vận tốc của chất điểm là $v(t)=s'(t)=-3t^2+12t=-3(t-2)^2+12\leq  12$.\\
	Dấu \rq\rq =\rq\rq\, xảy ra khi  và chỉ khi $ t=2 $.\\
	Vậy vận tốc chuyển động của  chất điểm đạt giá trị lớn nhất khi $t=2$.
	} 
\end{ex}

% nguồn: Ok_TraLoiNgan_Lop11_Lan1
\begin{ex}%[1D7V1-5]%[ma:25db80be]
	Cho chuyển động thẳng xác định bởi phương trình $S=-t^3+3 t^2+9 t$,  trong đó $t$ tính bằng giây và $S$ tính bằng mét. Tính vận tốc của chuyển động tại thời điểm gia tốc triệt tiêu.\\
	\shortans{12} 
	 \loigiai{
		Ta có: $v=S^{\prime}=-3 t^2+6 t+9$ và $a=v^{\prime}=-6 t+6$.\\ 
		Gia tốc triệt tiêu khi $a=0\Rightarrow  t=1$.\\
		Khi đó vận tốc của  chuyển động là $v(1)=12$ m/s.
	 }
\end{ex}

% nguồn: Ok_TraLoiNgan_Lop11_Lan2
\begin{ex}%[1D7V1-5]%[ma:c76dd2b2]
	Người ta xây dựng một cây cầu vượt giao thông hình parabol nối hai điểm có khoảng cách là $400$\,m. Độ dốc của mặt cầu không vượt quá $10^\circ$ (độ dốc tại một điểm được xác định bởi góc giữa phương tiếp xúc với mặt cầu và  phương ngang). Chiều cao giới hạn từ đỉnh cầu đến mặt đường là bao nhiêu mét? (làm tròn kết quả đến chũ số thập phân thứ nhất).
	\begin{center}
		\begin{tikzpicture}[line cap=butt,line join=miter,>=stealth,font=\normalsize]
				\path (-3.5,-49/48) coordinate (P) 
					(-2.5,-49/48) coordinate (Q) 
					(-2.5,-7/16) coordinate (R); 
				\tikzset{declare  function={xmin=-6; xmax=6; ymin=-3; ymax=0;
					f (\x)=-1/12*(\x)^2; 
					g  (\x)=7/12*(\x)+49/48;
					}, 
					smooth,samples=450 
				} 
				\draw (-6,-3) circle (1pt)  node[left]{$A$}--(6,-3) circle (1pt) node[right]{$B$}
					(-3.5,-49/48) circle (1pt) node[above left]{$P$}; 
				 \draw[dashed](0,0)--(0,-3) node[midway,right]{$h=?$}
					(-3.5,-49/48)--(-2.5,-49/48); 
				\clip (xmin,ymin) rectangle (xmax,ymax); 
				\draw  plot[domain=xmin:xmax] (\x,{f (\x)});
				\draw plot[domain=-4.5:-2.5] (\x,{g  (\x)});
				\draw pic[draw,,angle  radius=4mm]{angle=Q--P--R}
					pic[draw=none,"$10^\circ$",angle radius=12mm,font=\scriptsize]{angle=Q--P--R}; 
			\end{tikzpicture}
	\end{center}
	\shortans{17,6} 
	\loigiai{ 
		\begin{center}
			\begin{tikzpicture}[line cap=butt,line join=miter,>=stealth,font=\normalsize]
					 \path (-3.5,-49/48) coordinate (P)
						(-2.5,-49/48) coordinate (Q) 
						 (-2.5,-7/16) coordinate (R);
					 \tikzset{declare function={xmin=-7; xmax=7; ymin=-4; ymax=1.5;
						f (\x)=-1/12*(\x)^2; 
						g (\x)=7/12*(\x)+49/48; 
						}, 
						 smooth,samples=450
					} 
					\draw[dashed](-6,0)  node[above]{$-200$}|-(0,-3) node[below left]{$I$}
						(6,0)  node[above]{$200$}|-(0,-3)
						 (-3.5,-49/48)--(-2.5,-49/48);
					\draw (-3.5,-49/48) circle (1pt) node[above  left]{$P$}
						(-6,-3)  circle (1pt) node[left]{$A$}
						(6,-3)  circle (1pt) node[right]{$B$};
					\draw[->] (xmin,0)--(xmax,0) node[shift={(-100:7pt)}]{$  x $};
					\draw[->] (0,ymin)--(0,ymax) node[shift={(190:7pt)}]{$ y $}; 
					\fill (0,0) node[shift={(135:10pt)}]{$ O $}; 
					\clip (xmin,ymin) rectangle  (xmax,ymax);
					\draw  plot[domain=-6:6] (\x,{f (\x)});
					\draw  plot[domain=-4.5:-2.5] (\x,{g (\x)});
					\draw  pic[draw,,angle radius=4mm]{angle=Q--P--R}
						pic[draw=none,"$10^\circ$",angle  radius=12mm,font=\scriptsize]{angle=Q--P--R};
				\end{tikzpicture}
		\end{center}
		Chọn hệ trục toạ độ  như hình vẽ, sao cho đỉnh cầu là gốc tọa độ và mặt cắt của cây cầu có hình dạng parabol $(P)\colon y=-ax^2$ (với $a$ là hằng số dương).\\
		Giả sử điểm $P\in (P)$ có hoành độ  $x_0$, với $-200<x_0<200$.\\
		Hệ số góc của tiếp tuyến  của parabol bằng $k=y' \left(x_0 \right)=-2ax_0$.\\
		Hệ số góc xác định độ dốc của mặt cầu (độ dốc dương) là  $\left|k\right|=2a\left|x_0\right|\le 400a$.\\
		Vì độ dốc của  mặt cầu không vượt quá $10^\circ$ nên ta có
		\[400a\le \tan 10^\circ\Leftrightarrow a\le  \dfrac{\tan 10^\circ}{400}.
		 \]
		Chiều cao giới hạn từ đỉnh cầu  đến mặt đường là đoạn $OI$, cũng chính là độ lớn của tung độ điểm $B$ khi $a$ đạt giá trị lớn nhất.\\
		Do đó $OI=\left|-a\cdot 200^2\right|=17{,}6$\,m.\\ 
		Vậy chiều cao giới hạn  từ đỉnh cầu đến mặt đường là $17{,}6$\,m.
	 }
\end{ex}

% nguồn: Ok_TraLoiNgan_Lop11_Lan1
\begin{ex}%[1D7C1-5]%[1D5B3-3]%[ma:6cc63e53]
	\immini 
	 {
	Một vật gắn trên lò xo chuyển động theo phương ngang trên mặt phẳng nhẵn như hình bên. Phương trình chuyển động của vật được cho bởi $x=8\sin\left(\sqrt{2}\pi t+\dfrac{\pi}{3}\right)$, với $t$ tính bằng giây và $x$ tính bằng centimet. Tìm vận tốc và gia tốc của vật tại thời điểm $t=5$ giây (làm tròn đến chữ số  thập phân thứ nhất). Vật chuyển động theo hướng nào tại thời điểm đó?
	 }
	 {
	\begin{tikzpicture}[line cap=round,line join=round,>=triangle 45,x=1cm,y=1cm]
			 \draw [smooth,domain=3.13:15*pi,scale=0.1,samples=100] plot (\x,{4*sin (\x r)+7});
			 \fill [pattern = north east lines] (-0.2,0) rectangle (0,4);
			\draw[fill=gray]  (4.9,1.5)--(5.5,1.5)--(5.5,0)--(4.9,0)--cycle;
			\draw[->](0,0)--(6,0); 
			 \draw[->] (2.5,2.5)--(2.5,1.2);
			\draw (0,0) -- (0,4) 
				(2.5,0.1)--(2.5,-0.1) 
				(5.2,0.1)--(5.2,-0.1); 
			\draw (2.5,3) node {vị trí cân  bằng}
				(2.5,-0.3) node  {$O$}
				(5.2,-0.3) node  {$x$}
				(6,-0.3)  node {$x$}
				(3,-1) node  {Hình 9.1};
			 \draw (0,0.7)--(0.3,0.7)
				(4.72,0.7)--(4.9,0.7); 
		\end{tikzpicture}
	} 
	\shortans{151} 
	 \loigiai{
		Vận tốc và gia  tốc của vật tại thời điểm $t$ là
		\begin{itemize}
			 \item $v(t)=x'(t)=8\pi\sqrt{2}\cos\left(\sqrt{2}\pi t+\dfrac{\pi}{3}\right)$;
			 \item $a(t)=v'(t)=-16\pi^2\sin\left(\sqrt{2}\pi t+\dfrac{\pi}{3}\right)$.
		\end{itemize}
		Tại thời điểm $t=5$, vận tốc của vật là $v(5)=8\pi\sqrt{2}\cos\left(\sqrt{2}\pi \cdot 5+\dfrac{\pi}{3}\right)\approx -10{,}5$ cm/s.\\ 
		Gia tốc của vật  là $a(5)=-16\pi^2\sin\left(\sqrt{2}\pi \cdot 5+\dfrac{\pi}{3}\right)\approx 150{,}8$ cm/s$^2$.
	} 
\end{ex}


%%%%%%%%%%%%%%%% TỰ LUẬN %%%%%%%%%%%%%%%%

% nguồn: Ok_TuanLuan_Toan11_Lan1
\begin{bt}%[1D7H1-5]%[ma:763b5d32]
	Nhiệt độ cơ thể của một người trong thời gian bị bệnh được cho bởi công thức \break$T(t)=-0{,}1t^2+1{,}2t+98{,}6$, trong đó $T$ là nhiệt độ (tính theo đơn vị đo Fahrenheit) tại thời điểm $t$ (tính theo ngày). Tìm  tốc độ thay đổi nhiệt độ ở thời điểm $t=2$.
	 \loigiai{
		Đạo hàm của hàm số $T$ biểu  thị tốc độ thay đổi của nhiệt độ.\\
		 Ta có $T'(t)=-0{,}2t+1{,}2$.\\
		Vậy tốc  độ thay đổi nhiệt độ tại thời điểm $t=2$ là $T'(2)=-0{,}2\cdot 2+1{,}2=0{,}8^\circ F/$ngày.
	} 
\end{bt}

% nguồn: Ok_TuanLuan_Toan11_Lan2
\begin{ex}%[1D7H1-5]%[ma:f514852d]
	Nhiệt độ  cơ thể của một người trong thời gian bị bệnh được cho bởi công thức \break$T(t)=-0{,}1t^2+1{,}2t+98{,}6$, trong đó $T$ là nhiệt độ (tính theo đơn vị đo Fahrenheit) tại thời điểm $t$ (tính theo ngày). Tìm tốc độ thay đổi nhiệt độ ở thời điểm $t=2$.
	 \loigiai{
		Đạo hàm của hàm số $T$ biểu thị tốc  độ thay đổi của nhiệt độ.\\
		Ta có  $T'(t)=-0{,}2t+1{,}2$.\\
		Vậy tốc độ thay đổi nhiệt độ tại thời điểm  $t=2$ là $T'(2)=-0{,}2\cdot 2+1{,}2=0{,}8^\circ F/$ngày.
	} 
\end{ex}

% nguồn: Ok_TuanLuan_Toan11_Lan1
\begin{bt}%[1D7V1-5]%[ma:b72f0824]
	Một quả bóng được thả rơi tự do từ đài quan sát trên sân thượng  của toà nhà Landmark 81 (Thành phố Hồ Chí Minh) cao $461{,}3\,$ m xuống mặt đất, với phương trình chuyển động $s(t)=4{,}9t^2$. Tính vận tốc của quả bóng khi nó chạm đất, bỏ qua sức cản không khí (Đơn vị $m / s$, kết quả gần đúng làm tròn đến hàng phần chục).
	 \loigiai{
		Với bất kì $t_0$, ta  có
		$$s'\left(t_0\right)=\lim _{t \to t_0} \dfrac{s(t)-s\left(t_0\right)}{t-t_0}=\lim _{t \to  t_0} \dfrac{4{,}9t^2-4{,}9t_0^2}{t-t_0}=\lim _{t \to t_0} 4{,}9\left(t+t_0\right)=9{,}8t_0 .
		$$  Do đó, vận tốc của quả bóng tại thời điểm $t$ là $v(t)=s'(t)=9{,}8t$.\\
		Mặt  khác, vì chiều cao của toà tháp là $461{,}3\,$ m nên quả bóng sẽ chạm đất tại thời điểm $t_1$. \\
		Từ đó, ta có $4{,}9t_1^2=461{,}3  \Leftrightarrow t_1=\sqrt{\dfrac{461{,}3}{4{,}9}}$ (giây).\\
		Vậy vận tốc của quả bóng khi nó chạm  đất là
		$$v\left(t_1\right)=9{,}8t_1=9{,}8 \cdot \sqrt{\dfrac{461{,}3}{4{,}9}} \approx  95{,}1(\mathrm{~m} / \mathrm{s}).
		 $$
	 }
\end{bt}

% nguồn: Ok_TuanLuan_Toan11_Lan2
\begin{ex}%[1D7V1-5]%[ma:fe0a3998]
	Một quả bóng được thả rơi  tự do từ đài quan sát trên sân thượng của toà nhà Landmark 81 (Thành phố Hồ Chí Minh) cao $461{,}3\,$ m xuống mặt đất, với phương trình chuyển động $s(t)=4{,}9t^2$. Tính vận tốc của quả bóng khi nó chạm đất, bỏ qua sức cản không khí (Đơn vị $m / s$, kết quả gần đúng làm tròn đến hàng phần chục).
	 \loigiai{
		Với bất kì $t_0$, ta có 
		$$s'\left(t_0\right)=\lim _{t \to t_0} \dfrac{s(t)-s\left(t_0\right)}{t-t_0}=\lim _{t \to t_0} \dfrac{4{,}9t^2-4{,}9t_0^2}{t-t_0}=\lim _{t \to t_0}  4{,}9\left(t+t_0\right)=9{,}8t_0 .
		$$ Do đó, vận tốc của quả bóng tại  thời điểm $t$ là $v(t)=s'(t)=9{,}8t$.\\
		Mặt khác, vì chiều cao của toà tháp là $461{,}3\,$ m nên quả bóng sẽ chạm đất tại thời  điểm $t_1$. \\
		 Từ đó, ta có $4{,}9t_1^2=461{,}3 \Leftrightarrow t_1=\sqrt{\dfrac{461{,}3}{4{,}9}}$ (giây).\\
		Vậy vận tốc của quả bóng  khi nó chạm đất là
		$$v\left(t_1\right)=9{,}8t_1=9{,}8 \cdot \sqrt{\dfrac{461{,}3}{4{,}9}}  \approx 95{,}1(\mathrm{~m} / \mathrm{s}).
		$$ 
	} 
\end{ex}
